\chapter*{Prologue: why scattering still defines quantum dynamics}
% BEGIN CURATED SUBJECT INDEX
\index{spectrum!continuous}
\index{boundary condition!incoming and outgoing}
% END CURATED SUBJECT INDEX
\addcontentsline{toc}{chapter}{Prologue: why scattering still defines quantum dynamics}
\markboth{Prologue}{Why scattering still defines quantum dynamics}
\label{chap:prologue-scattering}

Quantum mechanics can be taught as finite-dimensional linear algebra.  That
language is powerful: qubits, spin systems, variational circuits, and
finite-basis calculations all live naturally in it.  Yet a finite matrix does
not by itself contain position or momentum as an unbounded observable, a
continuous spectrum, an infinite-volume limit, an asymptotic particle, or a
radiation condition.  Those are not decorative complications.  They are the
structures through which a dynamical theory meets an experiment.

Scattering provides the operational grammar.  One prepares a state in a
region where a comparison dynamics is understood, lets it traverse an
interaction region, and detects an outgoing state after another long stretch
of comparison dynamics.  The idealization involves three limits at once:
large separation, long time, and increasingly sharp control of the detector.
The mathematical theory exists to say when these limits define an invariant
map between preparation and detection data.  Resonance theory begins where
that map acquires a long-lived intermediate time scale.

\section*{A genealogy of the connection problem}
% BEGIN CURATED SUBJECT INDEX
\index{unbounded operator}
\index{topology}
\index{spectral theorem}
\index{spectral measure}
\index{spectrum!continuous}
\index{Sturm--Liouville theory}
\index{wave packet}
\index{boundary condition!incoming and outgoing}
\index{Jacobi coordinate}
% END CURATED SUBJECT INDEX
\phantomsection
\addcontentsline{toc}{section}{A genealogy of the connection problem}

The concepts assembled in this book were not invented in a single episode.
Hamilton--Jacobi theory encoded dynamics in a generating function.  The
Sturm--Liouville problem tied differential equations to orthogonality,
completeness, and boundary conditions.  Green functions made an inverse
operator into a propagation kernel, while Fourier analysis turned translation
into spectral multiplication.  Classical wave scattering supplied incoming
and outgoing radiation conditions.  Rutherford scattering then made the
large-distance deflection law a microscopic probe.  Quantum mechanics
inherited all of these ideas and forced them to coexist with probability and
operator domains.

The decisive mathematical step of the 1920s and 1930s was not merely to
replace observables by matrices.  It was to understand operators that are not
bounded and spectra that are not discrete.  Position, momentum, and the
Hamiltonian are specified by their domains as well as their differential
expressions.  A continuum state is represented by a spectral measure and a
wave packet, not by an extra normalizable eigenvector.  This is why the
spectral theorem, projection-valued measures, and the topology of operator
limits belong near the beginning of a modern account of quantum mechanics.

From this viewpoint, scattering theory is not a specialist's afterthought.
It is the part of quantum theory that keeps track of continuous spectrum,
infinite spatial extent, comparison dynamics, and the conversion of a
Hamiltonian into measured rates.  The same architecture reappears in nuclear
reactions, atomic ionization, condensed-matter transport, quantum field
theory, and black-hole perturbations.

\section*{The truncation contract}
% BEGIN CURATED SUBJECT INDEX
\index{operator!closed}
\index{Hilbert space}
\index{topology}
\index{spectrum!essential}
\index{resolvent}
\index{wave operator}
\index{Gauss law}
% END CURATED SUBJECT INDEX
\phantomsection
\addcontentsline{toc}{section}{The truncation contract}

Every computation is finite.  The important question is therefore not
whether one truncates, but what statement allows the finite problem to stand
for an infinite-dimensional one.  We shall call that statement the
\term{truncation contract}.  A complete contract names five ingredients:

\begin{enumerate}[label=(\roman*)]
  \item the parent Hilbert space and the closed operator or quadratic form;
  \item the projection, grid, basis, energy window, or channel selection;
  \item the topology in which the truncated objects are asserted to converge;
  \item the observable and error measure to be controlled; and
  \item the order in which basis, volume, time, regulator, and detector limits
  are taken.
\end{enumerate}

For an isolated bound state below the essential spectrum, the min--max
principle can provide such a contract.  For a scattering amplitude, one also
needs control of wave operators or resolvent boundary values.  For a
complex-scaled resonance, a small matrix residual is insufficient unless the
analytic-deformation hypotheses and stability window are checked.  For a
quantum simulation of a gauge theory, solving Gauss law can reduce the number
of qubits but can also reorganize locality and tensor factorization.  The
finite model becomes physically meaningful only when its contract states
which of these structures survives the reduction.

The finite-dimensional theory remains genuine dynamics: a unitary matrix can
generate nontrivial time evolution.  The limitation is more specific.  No
fixed finite dimension can encode, without a limiting prescription, the
unbounded observables and asymptotic structures on which scattering depends.
This book therefore uses finite matrices as controlled representations of a
declared parent problem, never as a replacement for declaring that problem.

\section*{Preparation, interaction, and detection}
% BEGIN CURATED SUBJECT INDEX
\index{Hilbert space}
\index{scattering matrix}
\index{boundary condition!incoming and outgoing}
\index{scattering!Coulomb}
\index{scattering!three-body Coulomb}
\index{distorted-wave approximation}
% END CURATED SUBJECT INDEX
\phantomsection
\addcontentsline{toc}{section}{Preparation, interaction, and detection}

Let $H_0$ denote the dynamics used to describe separated particles and let
$H$ contain the interaction.  A packet $\phi$ prepared in the distant past is
compared with the exact evolution through the strong limit
\begin{equation}
  \Omega^{(+)}\phi
  =\lim_{t\to-\infty}
  \rme^{\rmi Ht/\hbar}\rme^{-\rmi H_0t/\hbar}\phi,
  \label{eq:prologue-wave-operator}
\end{equation}
when that limit exists.  The analogous future limit defines
$\Omega^{(-)}$.  The scattering operator
\begin{equation}
  \Scattering=\Omega^{(-)\dagger}\Omega^{(+)}
  \label{eq:prologue-scattering-operator}
\end{equation}
maps incoming asymptotic data to outgoing asymptotic data.  The detailed
hypotheses and derivation are postponed to Stage II; here the formula fixes
the meaning of the symbol.  In particular, the intertwining relation implies
that $\Scattering$ commutes strongly with the comparison Hamiltonian $H_0$,
not with a generic interacting Hamiltonian.

The phrase ``free comparison'' is already conditional.  A short-range force
permits ordinary plane-wave asymptotics.  Coulomb scattering instead requires
Coulomb-distorted waves and a logarithmic phase.  With three charged bodies,
the simultaneous asymptotic regions are considerably more intricate, and no
single naive product of two-body plane waves solves the full problem.  In QED
the same long-range obstruction appears as soft radiation: a charged
asymptotic state generally does not belong to the undressed vacuum Fock
representation.  Thus the choice of asymptotic Hilbert space is part of the
physics, not merely a convention at the end of a calculation.

\section*{A diagnostic laboratory before the formalism}
% BEGIN CURATED SUBJECT INDEX
\index{operator!adjoint and self-adjoint}
\index{Wronskian}
\index{boundary condition!incoming and outgoing}
\index{scattering!Coulomb}
\index{continuum response}
% END CURATED SUBJECT INDEX
\phantomsection
\addcontentsline{toc}{section}{A diagnostic laboratory before the formalism}

Four elementary calculations will serve as recurring diagnostics.

\begin{description}[style=nextline]
  \item[A packet in a square well.]
  Separate the normalizable bound-state sum from the continuum integral and
  verify norm conservation.  A finite box discretizes the latter; the level
  spacing must disappear from observables as the box is enlarged.

  \item[A packet crossing a finite barrier.]
  Compute reflection and transmission from Wronskians.  Continue the incoming
  coefficient and locate its zeros.  This is the smallest model in which the
  same analytic datum appears as a scattering pole and an outgoing solution.

  \item[Hydrogenic motion.]
  Derive the self-adjoint radial problem, keep the Coulomb phase in its
  asymptotics, and distinguish the discrete spectrum from the continuum.  A
  plane-wave ansatz at infinity fails for a reason visible in the differential
  equation itself.

  \item[A finite-basis resonance calculation.]
  Vary basis size, nonlinear scale, and complex-rotation angle separately.
  A stable pole and a rotating discretized continuum play different roles;
  both are needed to reconstruct a response function.
\end{description}

These are not entrance examinations.  They are compact tests of whether a
representation still remembers its parent dynamics.  Each will be rebuilt in
later Stages with enough detail that the reader can reproduce the calculation.

\section*{The organizing claim of this book}
% BEGIN CURATED SUBJECT INDEX
\index{operator!adjoint and self-adjoint}
\index{Hilbert space}
\index{resolvent}
\index{Jost function}
\index{boundary condition!incoming and outgoing}
\index{exact WKB}
\index{complex scaling}
\index{Zel'dovich regularization}
\index{rigged Hilbert space}
\index{exceptional point}
% END CURATED SUBJECT INDEX
\phantomsection
\addcontentsline{toc}{section}{The organizing claim of this book}

The central claim is that a quantum resonance is best understood as a global
connection problem.  A self-adjoint Hamiltonian supplies real spectral and
probabilistic data.  Local solutions of its differential equation are
continued between asymptotic sectors.  The coefficient of the prohibited
incoming branch is an analytic function; its zero defines the resonance.
Exact WKB makes that construction nonperturbative and calculable.  Resolvent
continuation, Jost functions, Zel'dovich regularization, complex scaling, and
rigged Hilbert spaces are then different realizations of the same connection
datum, with different domains and topologies made explicit.

This order matters pedagogically.  We first define each mathematical object
where it is needed.  We then construct scattering and its analytic
continuation.  Only afterward do the Laboratories reuse those results in
radial motion, weakly bound nuclei, exceptional points, and black-hole
quasinormal modes.  References document provenance and extensions, but the
calculational route is contained in the book.
